\chapter{Kết quả kiểm tra và bàn giao}
\section{Kiểm tra lại revision bàn giao}
\noindent\begin{tabularx}{\linewidth}{@{}p{.37\linewidth}X@{}}
\toprule
\textbf{Kiểm tra} & \textbf{Kết quả quan sát}\\
\midrule
Pytest & 82 ca đạt trên Python 3.12. Bao gồm giữ chỗ, trả xe, hóa đơn, vòng đời sạc,
mô phỏng, trạng thái phiên và xử lý bản tin. Gói đã giải nén cũng chạy đủ 82 ca.\\
Kiểm tra kỹ thuật S2 & Đạt cấu trúc hồ sơ, truy vết 17 use case, cú pháp 54 sơ đồ PlantUML và kiểm thử.\\
Kiểm tra kỹ thuật S3 & Đạt chế độ \texttt{--technical-only} với 17 use case và 73 mục mô tả phương thức.\\
Kiểm tra trình duyệt & Hai cửa sổ SS3 đã mở trước khi phát và tự hiện cảnh báo.
SS7 từ chối preview tạo trước sự kiện.\\
\bottomrule
\end{tabularx}

\section{Bản tin trong chuỗi demo}
\begin{description}
\item[Mã] \texttt{b120ca85a0fe42c4932a3e9a362d6a14}
\item[Đích] HUB-001, cổng 1, loại \texttt{OFFLINE}
\item[Thời điểm phát] 10:23:50.930330 ngày 24/11/2026 (UTC+7)
\item[Thời điểm ghi SQLite] 10:23:50.930358 ngày 24/11/2026 (UTC+7)
\item[Quan sát trên UI] 2,27 và 3,36 giây kể từ lúc gọi tiến trình phát, theo thứ tự
kiểm tra cửa sổ. Khoảng đo gồm thời gian khởi động tiến trình và lấy kết quả UI.
\end{description}
Đối chiếu Hình~\ref{fig:demo-10-ss7-stale.png} và
\path{docs/report/submission_4/evidence/demo-run.json}.
Kết quả áp dụng cho lần chạy local này. Chỉ tiêu tải đồng thời, độ sẵn sàng và
độ trễ khi triển khai trên mạng chưa có phép đo tương ứng.

\section{Thành phần gói nộp}
Gói nộp gồm một PDF báo cáo cuối và thư mục \path{source/} chứa mã Python,
seed, tests, cấu hình chạy, nguồn LaTeX và dữ liệu liên quan được dùng khi dựng báo cáo.
Mã được chốt tại \texttt{aae0877}; tệp LaTeX và bằng chứng trình diễn bổ sung cho Submission 4.
Cách dựng báo cáo và cách tìm các nguồn tham chiếu nằm tại
\path{docs/report/submission_4/README.md}.
Tệp ZIP là gói bàn giao. PDF có thể dùng riêng nếu LMS cung cấp ô tải báo cáo.
Tại thời điểm chuẩn bị gói, ứng dụng đã được kiểm chứng local. Triển khai server
nhóm cần máy đích và quyền truy cập để kiểm tra từ máy khác.

\section{Nguồn đối chiếu}
\begin{enumerate}
\item Đề \path{BTL_SoftwareEngineering_HK261_v1.pdf}, trang 3--6, và yêu cầu bổ sung của
giảng viên về báo cáo, mã nguồn và giới hạn dung lượng ZIP.
\item Submission 1, báo cáo phân tích và \path{docs/requirements-analysis.md}.
\item Submission 2, nguồn thiết kế giao diện, tương tác và ma trận truy vết.
\item Submission 3, báo cáo thiết kế chung và tài liệu lớp/phương thức theo phân hệ.
\item Hướng dẫn và thiết kế hiện hành tại \path{docs/current/}; mã và test từ revision \texttt{aae0877}.
\end{enumerate}
